
\chapter{Discussion\label{discussion}}

\section{The \jlc\label{TheJetCorr}}

Various different measurements indicate that the \jlc is dominantly produced by vacuum fragmentation.  The spectra of particles in the \jlc shown in \Fref{RidgeVSJetVSPtAssoc} demonstrated behavior suggesting that the \jlc is dominated by fragmentation.  For instance, the spectra of particles in the \jlc have shallower slopes than the inclusive spectra and become shallower the higher the momentum of the trigger particle.  Likewise, the minimal dependence of the \jly on the angle relative to the reaction plane shown in \Fref{RxnPlaneDependence} and the agreement between the di-hadron fragmentation functions in \dAu and \Au shown in \Fref{zTRidgeSubtracted} also suggest that the \jlc is dominantly produced by fragmentation.  The lack of trigger particle dependence demonstrated in \Fref{JetEnergyTrigPtIdTrig}, \Fref{JetEnergyAssocPtIdTrig}, and \Fref{JetEnergyNpartIdTrig} is expected if the \jlc were dominated by fragmentation, however, strange particle production in jets is still not understood well enough to draw strong conclusions on this basis.

The particle composition of the \jlc further indicates that this feature is dominantly produced by fragmentation.  Particle ratios in \pp collisions at high \pT are likely to be dominated by fragmentation.  \Fref{RidgeParticleComposition} demonstrated that the baryon to meson ratio in the \jlc in central \Au collisions is comparable to that in \pp collisions for both strange and non-strange particles; \Fref{InclusiveRatioPlot} demonstrates that the baryon to meson ratio in the \jlc in \Cu is also comparable to the inclusive particle ratios in \pp collisions.

\begin{figure}
  \centering\resizebox{12cm}{!}{
  \includegraphics{/home/madscientistgirl/HarrisLab/Thesis/latex/results/InclusiveRatioPlot.eps}}\\

  \caption[Comparison of inclusive \baryontomeson to that in the \jlc and the \ridge in \Cu and \Au collisions at \sNNtwohundred]
  {Comparison of the inclusive \baryontomeson in \Cu \cite{Timmins:2007fe} and \Au \cite{Lamont:2004qy} collisions at \sNNtwohundred to that in the \jlc in 0-60\% central \Cu and 0-12\% central \Au collisions at \sNNtwohundred and to the \ridge in 0-12\% central collisions in \Au collisions at \sNNtwohundred. \cite{JanaWWND2007}.
  \label{InclusiveRatioPlot}}
\end{figure}

There is some indication that the \jlc may be slightly modified, particularly in central \Au.  \Fref{JetEnergyNpart} shows roughly a 50\% increase in \njet from \dAu collisions to central \Au collisions at \sNNtwohundred, and \Fref{JetEnergyWidths} indicates that the \jlc is slightly wider in \Au collisions at \sNNtwohundred.  The data are consistent with the same percentage increase in \Au collisions at \sNNsixtytwo, however, the error bars are too large to determine whether this increase is also present in collisions at \sNNsixtytwo.  There are three mechanisms which may lead to this increase.  First, both the \dphi and the \deta methods assume that the \ridge is independent of \deta, and if this is not true, some of the \ridge may be misidentified as part of the \jlc.  Collisions with higher \npart also have a larger \ridge and the \ridge increases with \npart.  Second, the slight difference in yields may also be due to a modification of the jet spectrum leading to the \jlc in \Au collisions.  For instance, if energy loss in the medium leads to a greater percentage energy loss for lower energy partons, the spectrum that contributes to the \jlc would be less steep in central \Au collisions.  This would mean a greater percentage of high energy partons would contribute to the \jlc.  Since higher energy partons fragment into more particles, the \jly would be greater.  Third, it is also possible that the \jlc is modified because of interactions of the parton with the medium and that this modification increases with the system size.  For instance, if some partons fragment in the medium, the fragmentation functions may be slightly modified in \AplusA collisions relative to \pp collisions.

The lack of dependence on the collision system of the \jlc demonstrated in the \pttrig dependence shown in \Fref{JetEnergyTrigPt}, the \ptassoc dependence shown in \Fref{JetEnergyAssocPt}, the \npart dependence shown in \Fref{JetEnergyNpart}, and the widths in \dphi and \deta shown in \Fref{JetEnergyWidths} further support the hypothesis that the \jlc is dominantly produced by fragmentation.  While each of these independent measurements has large errors, they all systematically indicate that the \jlc is comparable to what would be expected from fragmentation.



\subsection{Comparisons to \py}

In order to test whether these data are understood quantitatively, they are compared to \py \cite{Sjostrand:2003wg}.  \py is a Monte Carlo event generator for collisions such as \ee and \pp which incorporates analytical results and QCD-based models.  \py 6.4.10 was used with tune A which has been tuned to match data from the Tevatron at \sNN = 1.8 TeV \cite{Field:2002vt} and describes spectra of $\pi$ and p well at RHIC energies \cite{Adams:2006xb,Adams:2006nd,STARStrangePP200}.  To determine the \jly in \py, \pp collisions at \sNNtwohundred are simulated.  The analysis described in \Cref{CorrMethod} is repeated except that no separation of the \jlc and the \ridge is necessary because there is no \ridge in \pp.  The combinatorial background is assumed to be constant because there is no hydrodynamical flow in \py collisions.  Enough events were generated for the statistical error bars to be negligible.  Systematic error bars were not determined.

\begin{figure}
  \centering\resizebox{10cm}{!}{
  \includegraphics{/home/madscientistgirl/HarrisLab/Thesis/latex/results/TrigPtPythia6.eps}}\\
  \caption[\pttrig dependence of \njet for \sNNsixtytwo and \sNNtwohundred compared to \py]
  {\pttrig dependence of \njet for \stdassoc for 0-60\% central \Cu and 0-80\% central \Au collisions at \sNNsixtytwo and minimum bias \dAu, 0-60\% central \Cu and 0-12\% central \Au collisions at \sNNtwohundred.  Data from \dAu and \Au collisions at \sNNtwohundred are from \cite{JanaFrance,JanaWWND2007,JanaQM06,JanaHP06,JanaPANIC05,JanaPrivComm}.  Lines indicate predictions from \py.
  \label{JetEnergyTrigPtPY}}
\end{figure}

\Fref{JetEnergyTrigPtPY} shows the dependence of the \jlc on \pttrig for all systems and energies studied, as shown in \Fref{JetEnergyTrigPt}, compared to predictions from \py.  Since the \dAu data from di-hadron correlations match \pp where comparisons are possible, \py should be able to describe the \dAu data.  However, since there are no differences observed between \dAu, \Cu, and \Au, \py should be able to describe all of the data.  \py overestimates the data at low \pT but matches the data at higher \pT.  The data used to tune \py were from considerably higher jet energies and therefore may not constrain this kinematic region well.  Deviations of \py from the data could be explained if either \py does not use the correct fragmentation functions or if the jet spectrum which produces the \jlc is not correct in \py.  Preliminary studies of the fragmentation functions in heavy ion collisions indicate that \py describes the fragmentation functions in heavy ion collisions \cite{jetrecoJoernHP}, although the jet sample used for these studies may be biased towards unmodified jets and the average jet energy which leads to the \jlc presented in \Fref{JetEnergyTrigPtPY} is considerably lower.

\begin{figure}
  \centering\resizebox{10cm}{!}{
  \includegraphics{/home/madscientistgirl/HarrisLab/Thesis/latex/results/AssocPtPythia6.eps}}\\

  \caption[\ptassoc dependence of \njet for \sNNsixtytwo and \sNNtwohundred compared to \py]{\ptassoc dependence of \njet for \stdassoc for \stdtrig for 0-60\% central \Cu and 0-80\% central \Au collisions at \sNNsixtytwo and minimum bias \dAu, 0-60\% central \Cu and 0-12\% central \Au collisions at \sNNtwohundred.  Data from \dAu and \Au collisions at \sNNtwohundred are from \cite{JanaFrance,JanaWWND2007,JanaQM06,JanaHP06,JanaPANIC05,JanaPrivComm}.  The lines that match the data indicate fits, as shown in \Fref{JetEnergyAssocPtPy}, and the thicker, shaded lines indicate predictions from \py.
  \label{JetEnergyAssocPtPy}}
\end{figure}

% \begin{table}
% \begin{center}%The inverse slope parameter from a fit to $\pi^-$ as a function of \pT in \Au from \cite{Abelev:2007ra} above 1.0 \GeV is $k$ = 280.9 $\pm$ 0.4 \MeV for \sNNsixtytwo and is $k$ = 330.9 $\pm$ 0.3 \MeV for \sNNtwohundred.
% \caption{Inverse slope parameter k (\MeV) of \ptassoc for fits of data in \Fref{JetEnergyAssocPt} and \Fref{JetEnergyAssocPtPy} to A$e^{-p_{T}/k}$ compared to the inclusive spectra and \py.  Statistical errors only.
% }\label{EnergyDependenceSlopeTablePy}
% \begin{tabular}[b]{c|c|c}
% \hline
%  & \sNNsixtytwo & \sNNtwohundred\\
%    & k (\MeV) & k (\MeV) \\ \hline
% %Au \ridge&  & 438 $\pm$ 4   \\ 
% \Au & 291 $\pm$ 28 & 478 $\pm$ 8 \\
% \Cu & 359 $\pm$ 41 & 424 $\pm$ 20 \\ 
% \dAu &  & 469 $\pm$ 8 \\ \hline
% \Au \pion spectra & 280.9 $\pm$ 0.4 & 330.9 $\pm$ 0.3 \\
% \Cu \pion spectra & 277 $\pm$ 11 & 304 $\pm$ 13\\
% \dAu inclusive spectra &  & 350 $\pm$ 5\\ \hline
% \py & 171 $\pm$ 1 & 250.0 $\pm$ 0.9\\ \hline%PY6.4.1 10 June 09
% \end{tabular}
% \end{center}
% \end{table}

\Fref{JetEnergyAssocPtPy} shows the dependence of the \jly on \ptassoc, as shown in \Fref{JetEnergyAssocPt}, compared to \py.
% The inverse slope parameters from the inclusive spectra and \py from fits are compared  in \Tref{EnergyDependenceSlopeTablePy} to the inverse slope parameters from fits of the \jlc given in \Tref{EnergyDependenceSlopeTable}.  Fits to the inclusive \pion spectra from 1.0 \GeV $<$ \pT $<$ 3.0 are given using \Au data from \cite{Abelev:2007ra}, \Cu data from \cite{PHENIX:2008cx}, and \dAu data from \cite{PHENIX:2007by}.  The inverse slope parameter of the \jlc is systematically higher than that of the inclusive spectra.  
\py does not correctly describe the spectrum of particles in the \jlc in collisions.

\begin{figure}
  \centering\resizebox{10cm}{!}{
  \includegraphics{/home/madscientistgirl/HarrisLab/Thesis/latex/results/JetNpartPythia6.eps}}\\

  \caption[\npart dependence of \njet for \sNNsixtytwo and \sNNtwohundred compared to \py]
  {\npart dependence of \njet for \stdtrig and \stdassoc for \Cu and \Au collisions at \sNNsixtytwo and \dAu, \Cu and \Au collisions at \sNNtwohundred.  Data from \dAu and \Au collisions at \sNNtwohundred are from \cite{JanaFrance,JanaWWND2007,JanaQM06,JanaHP06,JanaPANIC05,JanaPrivComm}.  Lines indicate predictions from \py.
  \label{JetEnergyNpartPy}}
\end{figure}

\begin{figure}
  \centering\resizebox{15cm}{!}{
  \includegraphics{/home/madscientistgirl/HarrisLab/Thesis/latex/results/WidthsPythia6.eps}
}\\

  \caption[Collision energy and system size dependence of widths compared to \py]
  {Collision energy and system size dependence of Gaussian widths in \dphi and \deta as a function of \pttrig, \ptassoc, and \npart  for \Cu and \Au collisions at \sNNsixtytwo and \dAu, \Cu and \Au collisions at \sNNtwohundred.  Kinematic cuts and centralities are the same as used in \Fref{JetEnergyTrigPt}, \Fref{JetEnergyAssocPt}, and \Fref{JetEnergyNpart} except that the \Au data at \sNNtwohundred for the \pttrig and \ptassoc are from 0-80\% central collisions.  Lines indicate predictions from \py.
  \label{JetEnergyWidthsPy}}
\end{figure}

The \jly expected in \py is compared to the data in \Fref{JetEnergyNpart} in \Fref{JetEnergyNpartPy}.  \py describes the data in this kinematic range well.  The widths in \dphi and \deta shown in \Fref{JetEnergyWidths} are compared to \py in \Fref{JetEnergyWidthsPy}.  \py underestimates the widths in \dphi, particularly for the lowest \pttrig.  While the \dAu data are higher than \py, the widths in \AplusA are larger than \dAu and the width grows with \npart.  The widths in \deta generally agree with the data except for the lowest \ptassoc in \Cu collisions and for the lowest \pttrig in \Au collisions at \sNNtwohundred.  The width in \deta in \Au collisions at \sNNtwohundred matches \py in the most peripheral collisions but by the most central collisions the width in \deta has almost doubled.%  \py describes the widths reasonably well, although it genderally underestimates the widths in \dphi for the lowest \pttrig and to overestimate the width in \deta for the lowest \ptassoc.

Before these studies were done, \py was not necessarily expected to compare well with any aspect of di-hadron correlations because it is clear from the observation of jet quenching shown in \Fref{IconicJetQuench} and the \ridge that di-jets are strongly modified in \AplusA collisions.  Additionally, \py has not been tuned to data from di-hadron correlations.  \py not only qualitatively describes the \jlc but also describes the \jlc fairly well quantitatively.  The fact that few differences are observed between the \jlc in \dAu, \Cu, and \Au collisions implies that there are no strong nuclear effects modifying the \jlc in \AplusA collisions.  This means that the deviations of \py from the data at the lowest \pttrig and \ptassoc can be interpreted as a need for better tuning of \py to lower \pT data.  Deviations from \pp, peripheral \AplusA, and \py observed in the central \Au can be understood to result from either modifications of the \jlc or evidence of a slight dependence of the \ridge on \deta.

The agreement of the \jlc with \py means that a leading order pQCD model such as \py can be used to understand how the kinematic cuts applied in di-hadron correlations affect the jet sample.  This demonstrates that di-hadron correlations are an effective means for studying jets in heavy ion collisions and that the \jlc can be calculated from first principles.  Therefore, the \jlc provides a means to determine the distribution of jet energies of jets leading to di-hadron correlations, a useful tool for studies of both the \as and the \ridge.  It also means that models for the \ridge which would affect the per trigger yield of the \jlc are not correct.

\section{The \ridge}

\subsection{Theoretical descriptions of the \ridge}

The \ridge is a novel feature only observed in heavy ion collisions.  Multiple different models have been proposed to explain the \ridge.  These models build on the theoretical frameworks used to describe heavy ion collisions discussed in \Cref{introduction}.

%Many of these models incorporate hydrodynamical phenomena.  Others link the \ridge to initial condititions in a Color Glass Condensate or a plasma instability.

\subsubsection{Radial flow plus trigger bias}

\begin{figure}
  \centering\resizebox{6cm}{!}{\rotatebox{0}{
  \includegraphics{/home/madscientistgirl/HarrisLab/Thesis/latex/theory/radialflowtriggerbias.eps}}}\\

  \caption[Schematic diagram depicting the radial flow plus trigger bias model]
  {Schematic diagram depicting the radial flow plus trigger bias model
  \label{RadialFlowSchematic}}
\end{figure}

The radial flow plus trigger bias model explains the \ridge as an incidental rather than a causal correlation between particles \cite{Voloshin:2009fd}.  As discussed in \Cref{introduction}, there is evidence that the medium produced in central heavy ion collisions at RHIC can be described as a fluid and that there are pressure gradients in this fluid.  Hydrodynamical models predict the expansion of the medium through radial flow, meaning that particles in the bulk get a radial velocity boost so that they are mostly emitted perpendicular to the surface of the medium.  In the radial flow plus trigger bias model, hard partons are also emitted mostly perpendicular to the surface of the medium.  The hard partons may either get a radial velocity boost from radial flow, or they may be surface-biased jets.  Either mechanism means that both the particles from the jet and particles from the medium come mostly from the surface, resulting in a correlation between these particles.  This is depicted schematically in \Fref{RadialFlowSchematic}.  The \ridge is then made up of particles in the bulk that are correlated with the hard parton through this mechanism.  Because the distribution of particles in the bulk is broad in pseudorapidity, the \ridge is then broad in pseudorapidity.


\begin{figure}
  \centering\resizebox{7cm}{!}{\rotatebox{0}{
  \includegraphics{/home/madscientistgirl/HarrisLab/Thesis/latex/theory/SergeiBeta0.4trig3assoc1_2.eps}}}\\

  \caption[Di-hadron correlation predicted in the radial flow plus trigger bias model]
  {Di-hadron correlation predicted in the radial flow plus trigger bias model for a radial velocity boost $\beta_{r}$ = 0.4, and \trigrange{3.0}{20.0}, \assocrange{1.0}{2.0} \cite{Pruneau:2007ua}
  \label{SergeiRidge}}
\end{figure}

The calculation in \cite{Pruneau:2007ua} by Pruneau, Gavin, and Voloshin (PGV) uses \py for generating \pp events and gives the particles in these events a radial velocity boost to mimic the effect of radial flow.  This calculation does not conserve momentum but is able to recreate many features observed in the data.  It produces an enhancement on the \ns similar to the \ridge and a dip on the \as similar to the dip observed in the data.  \Fref{SergeiRidge} shows the di-hadron correlation produced in this model for similar kinematic cuts to those used in the data in this thesis and with a radial velocity boost comparable to that inferred from comparisons of the inclusive particle spectra in central \Au collisions at \sNNtwohundred to a Blast Wave model.  In this kinematic region, the \as is negligible.  No quantitative comparisons between this model and the data presented in this thesis have been done.

\begin{figure}
  \centering\resizebox{9cm}{!}{\rotatebox{0}{
  \includegraphics{/home/madscientistgirl/HarrisLab/Thesis/latex/theory/JunRidge.eps}}}\\

  \caption[Di-hadron correlation predicted in the radial flow plus trigger bias model]
  {Di-hadron correlation predicted for \trigrangevar{2.5} and \assocrangevartwo{1.0} in \cite{Takahashi:2009na}.
  \label{JunRidge}}
\end{figure}

Theoretical calculations involving both jets and hydrodynamics are difficult because hydrodynamics inherently describes processes that occur in a fluid at equilibrium and jets are inherently not at equilibrium.  Attempts to calculate the production of the \ridge through the radial flow plus trigger bias mechanism using a hydrodynamical framework have been done, although they do not yet agree with the data.  This may be due to the inherent difficulty in such calculations.  The calculation by Shuryak in \cite{Shuryak:2007fu} gets the width of the \ridge in azimuth a factor of two larger than observed in the data.  The calculation in \cite{Takahashi:2009na} by Takahashi et al is a full 3D hydrodynamical calculation which reproduces many of the general features in the data, as shown in \Fref{JunRidge}.  In this calculation a hard parton is mimicked by introducing hot cells into the hydrodynamical calculation.  In addition to producing a \ridge, this leaves a correlation which looks like the \jlc, even though there are no jets generated.

\subsubsection{Glasma flux tubes}
In the Color Glass Condensate (CGC) picture, the incoming nuclei are described as two sheets of Glasma.  When these sheets collide, strong color magnetic and electric fields form and in a short period of time these fields change from transverse to longitudinal \cite{McLerran:2008uj}.  In the CGC picture the \ridge is formed by flux tubes of longitudinal color electric and magnetic fields that form in this initial state \cite{Dumitru:2008wn,Gelis:2008sz}.  These correlations are enhanced by the radial flow plus trigger bias mechanism; without enhancement by radial flow the \ridge produced in this model is not large enough to be consistent with the data.

\begin{figure}
  \centering\resizebox{8cm}{!}{\rotatebox{-0}{
  \includegraphics{/home/madscientistgirl/HarrisLab/Thesis/latex/theory/QCDFieldsRidge2.eps}}}\\

  \caption[The \ridge produced in a Color Glass Condensate]
  {The \ridge produced in the CGC \cite{MizukawaQCDMagFields}.
  \label{QCDFieldsRidge}}
\end{figure}

Most quantitative comparisons with the data have been with the untriggered, or ``Soft'', \ridge \cite{Gavin:2008ev}, however, the triggered, or ``Hard'', \ridge should be described in this model.  Calculations of the ``Hard'' \ridge in this model are difficult.  The calculation of the ``Hard'' \ridge done in \cite{MizukawaQCDMagFields}, shown in \Fref{QCDFieldsRidge}, qualitatively reproduces the shape of the \ridge, however, the size of the \ridge relative to the \jlc is considerably smaller than what is observed in the data.  In this model the \jlc would not be affected because the parton would fragment after having left the medium.

\subsubsection{Recombination and the Correlated Emission Model}
Recombination was introduced in response to a baryon enhancement in the inclusive data.  In Recombination, hadrons form by the combination of partons in the medium.  The partons may either be thermal partons, those coming from the medium, or shower partons, those coming from hard scattering.  In this model the \jlc comes from shower partons or combinations of thermal and shower partons.  The inclusive spectra were fit to determine the thermal and shower contributions to the spectra and based on these fits it was predicted by Hwa that there would be no correlation observed for \myomega triggered correlations \cite{Hwa:2006vb}.  The data in \cite{bettysthesis,Abelev:2007zp} contradicted this prediction.  In \cite{Chiu:2007zr} Hwa explained the discrepancy by hypothesizing that the \ridge comes from ``phantom jets'', that the \myomega trigger is actually from the \ridge, and that the \myomega triggered correlation in central \Au collisions at \sNNtwohundred presented in \cite{bettysthesis,Abelev:2007zp} comprises only a \ridge but no \jlc.  There are not enough statistics to separate the \jlc from the \ridge in the \myomega triggered correlation.  This model presumes, however, that there is a jet which created the \ridge; the \myomega trigger is in the \ridge for this jet.

This model was used to calculate the \ridge in central \Au collisions at \sNNtwohundred and is able to describe the data \cite{Hwa:2009bh}.  No predictions for \Cu collisions or to collisions at \sNNsixtytwo have been made yet.  Hwa's Correlated Emission Model uses recombination to generate particles from the bulk and relies on the radial flow plus trigger bias model mechanism to get the correlated emission of particles from the bulk and from hard partons responsible for the \ridge \cite{Chiu:2008ht}.  Since some of the particles in the \jlc would come from a combination of thermal and shower partons, this model may lead to a modification of the \jlc.

\subsubsection{Momentum Kick Model}
Wong's Momentum Kick Model \cite{WongMomKick4} posits that the \ridge arises from collisions of hard partons with partons from the medium before fragmentation.  The medium partons are assumed to have a roughly uniform distribution in azimuth and pseudorapidity before the collision with the hard parton.  Collisions of the hard parton give medium partons momentum kicks in the direction of the hard parton, creating an azimuthal correlation between medium partons and the hard parton.  The momentum kick does not create a strong correlation in pseudorapidity because jets are dominantly produced at midrapidity \cite{WongMomKick4,Cheukyin}.  This scenario is depicted in \Fref{WongMomKickSchematic}.

\begin{figure}
  \centering\resizebox{6cm}{!}{
  \includegraphics{/home/madscientistgirl/HarrisLab/Thesis/latex/theory/momentumkickschematic.eps}}\\

  \caption[Schematic diagram showing the mechanism for momentum kick model.]
  {Schematic diagram exhibiting the mechanism for \ridge production in the momentum kick model \cite{WongMomKickMain}.  The solid red lines show the trajectory of the parton as it traverses the medium and the dashed red lines show the directions of the momenta of the hadrons from the fragmentation of the hard parton.
  \label{WongMomKickSchematic}}
\end{figure}

The Momentum Kick Model does not explain how the initial uniformity of the parton distribution at mid-rapidity arises.  It may come from string fragmentation, or from another mechanism such as a Color Glass Condensate initial state.  The Momentum Kick Model fits the experimental data on the \ridge and determines the initial parton momentum from these data.  It successfully fits the \ptassoc and \npart dependence of the \ridge \cite{WongMomKickMain}.  It also successfully predicted the shape of the \ridge at large \deta in PHOBOS using the data available at mid-rapidity from STAR \cite{Wenger:2008ts}.


\begin{figure}
  \centering\resizebox{7cm}{!}{
  \includegraphics{/home/madscientistgirl/HarrisLab/Thesis/latex/theory/MomentumKickEnergyDependence.eps}}\\

  \caption[Collision energy dependence in the Momentum Kick model]
  {Collision energy dependence of the \ridge in the Momentum Kick model \cite{WongMomKickMain} for the unidentified hadron correlations shown in \Fref{RidgeEnergyNpart}.
  \label{WongMomEnergyDependence}}
\end{figure}

The Momentum Kick Model has been used to describe the data on the collision energy dependence of the \ridge presented in this thesis.  This was not done by fitting the data from \sNNsixtytwo but by assuming that the size of the \ridge scales with the density of partons and that the density of partons scales with the density of participants \cite{WongMomKickMain}.  This energy dependence is shown in \Fref{WongMomEnergyDependence} and agrees with the data.


\begin{figure}
  \centering\resizebox{7cm}{!}{
  \includegraphics{/home/madscientistgirl/HarrisLab/Thesis/latex/theory/momkicketadist.eps}}\\

  \caption[Rapidity dependence of the \ridge yield in the Momentum Kick Model]
  {Rapidity dependence of the \ridge yield in the Momentum Kick Model \cite{WongMomKickMain}.
  \label{WongMomRapDep}}
\end{figure}

The Momentum Kick Model predicts a sharp drop in the \ridge as a function of the rapidity of the associated particle and \ptassoc spectra which are strongly dependent on rapidity.  These features are dependent on the kinematic cuts used in the analysis.  This is depicted in \Fref{WongMomRapDep}.  Data which extend beyond this region would be useful for testing the Momentum Kick Model.  The PHOBOS measurement \cite{Wenger:2008ts} extends to \deta = 4, however, there is no measurement of \ptassoc and the lower limit on \ptassoc is determined by the acceptance of the detector.  This lower limit is around 35 \MeV so this measurement cannot exclude the Momentum Kick Model.  STAR's acceptance is limited in rapidity, although higher statistics data sets may allow a high enough \ptassoc to test the prediction in \Fref{WongMomRapDep}.


\subsubsection{Bremsstrahlung gluons broadened by medium interactions}% needs better name
\begin{figure}
  \centering\resizebox{8cm}{!}{\rotatebox{-90}{
  \includegraphics{/home/madscientistgirl/HarrisLab/Thesis/latex/theory/JetsInMediumWind.eps}}}\\

  \caption[\ridge production through radiated gluons broadened by transverse flow]
  {\ridge production through radiated gluons broadened by transverse flow \cite{Armesto:2004pt}
  \label{JetsInMediumWind}}
\end{figure}
The hard partons which fragment into jets may emit bremsstrahlung gluons.  It is hypothesized in \cite{Armesto:2004pt} that the \ridge is the result of gluon bremsstrahlung in the medium broadened by transverse flow, as depicted schematically in \Fref{JetsInMediumWind}.  This model describes moderate broadening of the \jlc in \deta, however, it has difficulty with producing a structure as broad in \deta as the \ridge.  The calculations in \cite{Armesto:2004pt} predict a \ridge which is only about 10\% broader in \deta than the \jlc.

\begin{figure}
  \centering\resizebox{6cm}{!}{\rotatebox{0}{
  \includegraphics{/home/madscientistgirl/HarrisLab/Thesis/latex/theory/plasmainstability.eps}}}\\

  \caption[Broadening of the \jlc from plasma instabilities]
  {Broadening of the \jlc from plasma instabilities for two different jet energies \cite{Majumder:2006wi}.
  \label{PlasmaInstabilities}}
\end{figure}

The formation of the \ridge from bremsstrahlung gluons is also hypothesized in \cite{Majumder:2006wi}, however, in this model the gluons are broadened by strong, turbulent color fields.  These strong color fields are predicted to be produced in response to plasma instabilities.  If such fields exist, they would deflect brehmstralung gluons.  \Fref{PlasmaInstabilities} shows the broadening predicted in \cite{Majumder:2006wi} for two different jet energies; the structure in these figures is not broad enough to be consistent with the data.

\subsection{Discussion}
Data on the \ridge is extensive and there are several observations which may be key to distinguishing production mechanisms for the \ridge:
\begin{enumerate}
 \item  \textbf{The \jlc is not substantially modified.}  \Sref{TheJetCorr} discusses the overwhelming evidence that the \jlc is produced dominantly by vacuum fragmentation.  This means that the production of the \ridge must not substantially modify the \jlc.\label{JetNotMod}
 \item  \textbf{The particle ratios in the \ridge are comparable to those in the bulk.}  \Fref{RidgeParticleComposition} and \Fref{InclusiveRatioPlot} demonstrate that the \ridge has a particle composition similar to the bulk.  By contrast, the \jlc has a particle composition similar to what would be expected from vacuum fragmentation.\label{PartRatios}
 \item  \textbf{The \ridge is larger in collisions at \sNNtwohundred than at \sNNsixtytwo.}  This is shown in \Fref{RidgeEnergyNpart}.\label{EnergyDep}
 \item  \textbf{The \ridge is larger in plane than out of plane.}  This is shown in \Fref{RxnPlaneDependence}.  This corresponds to a larger \ridge when the path length traveled by a parton through the medium would be lower on average.\label{RxnDep}
 \item  \textbf{There is either no dependence on the trigger particle or a larger \ridge for a heavier trigger particle.}  \Fref{RidgeNpartIDTrig} shows the dependence of the \ridge on \npart for different trigger particle species.  It is unclear how correlated the systematic errors due to \vtwo are so it is difficult to conclusively say that there is mass dependence.  However, if the systematic errors are correlated, then the \ridge is larger for \lam trigger particles than for unidentified trigger particles and it is smaller for \kaon trigger particles than for an unidentified trigger particles.  \Fref{RidgeParticleComposition} shows that in the kinematic range of the trigger particle there are roughly as many protons as pions.  Since protons and pions are the most abundant particles, this means that the unidentified hadron trigger particles are roughly 50\% protons and roughly 50\% pions, making the average mass for unidentified triggers between the \lam mass and the \kaon mass.  This means that if the systematic errors in \Fref{RidgeNpartIDTrig} are correlated, the \ridge is larger for heavier trigger particles.  A smaller \ridge for heavier trigger particles is strongly disfavored by the data.\label{TrigDep}
\item  \textbf{The \ridge is broad in \deta.} The \ridge is roughly independent of \deta within the STAR acceptance.  There is some room for a slight \deta dependence, but there cannot be a strong dependence on \deta. The PHOBOS measurement \cite{Wenger:2008ts} demonstrated that the \ridge extends to \deta = 4, however, the associated particles were determined from single hits, so the kinematic reach was limited by the acceptance of the detector.  This limit is about \pT $>$ 35 \MeV.  The extent of the \ridge may cause some problems for causal models because a parton at midrapidity has to affect particles separated by a large gap in \deta.\label{Extent}
\end{enumerate}

 \subsubsection{Radial flow plus trigger bias mechanism and related models}
The Glasma Flux Tube model and the Correlated Emission Model both primarily produce the \ridge using the radial flow plus trigger bias mechanism.  The radial flow plus trigger bias mechanism naturally leads to agreement with many of the experimental observations.  Because hard partons fragment as in a vacuum in the PGV and Shuryak calculations, these models do not modify the \jlc, in agreement with \oref{JetNotMod}.  The Recombination model by Hwa may lead to some modification of the \jlc because some of the \jlc could come from recombination of shower and thermal partons rather than vacuum fragmentation.  A calculation needs to be done in this framework to ensure that the Recombination model is consistent with \oref{JetNotMod}.  In all models, particles in the \ridge come from the bulk, in agreement with \oref{PartRatios}.  Assuming that the reaction plane dependence calculated in the Correlated Emission Model \cite{Chiu:2008ht} is a common feature of the radial flow plus trigger bias mechanism, these models are also in agreement with \oref{RxnDep}.  The \ridge would also be expected to have roughly the same extent in \deta as inclusive particles, in agreement with \oref{Extent}.

It is not clear whether the energy dependence (\oref{EnergyDep}) is consistent with this model but the energy dependence would be readily calculated in the same way as done by PGV \cite{Pruneau:2007ua}.  The amount of radial flow can be determined by comparisons with the inclusive spectra.  As shown in \Fref{LowpTv2} for \vtwo, lighter particles are affected more by hydrodynamical flow.  If these models are dependent on a velocity boost for the surface bias of the hard parton, then a dependence on the mass of the trigger particle should be observed and lighter trigger particles should correspond to a larger \ridge.  This contradicts \oref{TrigDep}.  This means that if the radial flow plus trigger bias mechanism is the proper mechanism for producing the \ridge, the surface bias of the hard partons likely does not come from radial flow but from jet quenching.

 \subsubsection{The Momentum Kick Model}
In the Momentum Kick Model the \jlc is produced by vacuum fragmentation of the parton after it interacts with the medium.  The \jlc is not affected, in agreement with \oref{JetNotMod}.  The \ridge is created by medium partons so naturally has the same composition as the bulk, in agreement with \oref{PartRatios}.  As shown in \Fref{WongMomEnergyDependence}, the Momentum Kick Model describes the energy dependence of the \ridge, in agreement with \oref{EnergyDep}.  No calculations have been made for the dependence of the \ridge on trigger particle type, so it is unclear if this model is consistent with \oref{TrigDep}.

Naively the larger path length out of plane would lead to a larger \ridge, in contradiction to \oref{RxnDep}.  It may be possible to reconcile this with the data if the parton loses so much energy when it travels out of plane that it is not observed.  In that case only jets from the surface would be observed out of plane.  The reaction plane dependence of the \ridge in the Momentum Kick Model has not been calculated yet.

The Momentum Kick Model is the only causal model that produces a \ridge broad enough to be consistent with the data, in agreement with \oref{Extent}.  However, this model also has a sharp drop just outside of reach of current measurements.  This unusual feature, if observed, would be one signature that the Momentum Kick Model is the correct mechanism for the production of the \ridge.

\subsubsection{Gluon Bremsstrahlung Models}
In gluon bremsstrahlung models, the hard parton fragments in the vacuum after having emitted bremsstrahlung gluons, in agreement with \oref{JetNotMod}.  In these models, the \ridge would come from the fragmentation of gluons, which would naively seem to contradict \oref{PartRatios}, that the \ridge composition is similar to the bulk.  If the bremsstrahlung gluons are sufficiently soft that they reinteract with the medium, this may be able to be reconciled with the data.  No calculation of the collision energy dependence has been done, however, the expectations from these models seem to be consistent with \oref{EnergyDep}.  These models seem to have the same contradiction with \oref{RxnDep} that the Momentum Kick Model does, that naively the larger path length out of plane would lead to a larger \ridge.  If heavier trigger particles correspond on average to heavier partons, then lighter trigger particles should produce a larger \ridge because the amount of bremsstrahlung is proportional to $m^{-4}$.  This is in contradiction to \oref{TrigDep}.  These models also have difficulty producing a \ridge broad enough in pseudorapidity to match the data, in contradiction to \oref{Extent}.

% 
% The models for the \ridge can be divided into those where the \ridge is caused by interactions with the parton which creates the \jlc and those where the \ridge is correlated with the trigger particle through other mechanisms but is not caused by an interaction with the parton.  Causal models include the Momentum Kick Model, and models for gluon brehmstralung broadened by interactions with the medium.  Models for gluon brehmstralung have had difficulty producing a \ridge which is broad enough in pseudorapidity to match the data.  Causal models may have problems creating a \ridge which is both wide enough in pseudorapidity and as large as observed in the data.  To date gluon bremsstrahlung models have not been able to quantitatively describe the data; this does not mean that partons do not lose energy from these mechanisms, only that they are not likely to create the \ridge.  Indeed, these mechanisms may be responsible for the broadening of the \jlc in central \Au collisions at \sNNtwohundred.
% 
% Non-causal models include the radial flow plus trigger bias model, the Correlated Emission Model, and the Glasma flux tube model.  The radial flow plus trigger bias model qualitatively describes the di-hadron correlation in central \Au collisions at \sNNtwohundred.  Both the Correlated Emission Model and the Glasma flux tube model use the radial flow plus trigger bias mechanism.  All of the calculations for the \ridge so far in these models have either been limited to central \Au collisions at \sNNtwohundred or have attempted to describe untriggered di-hadron correlations.  The energy dependence should be calculable in a model such as \cite{Pruneau:2007ua} because the amount of radial flow is constrained by inclusive particle spectra.  Such a model should also be able to predict the dependence on \pttrig.  The calculation in \cite{Chiu:2008ht} indicates that the reaction plane dependence of the \ridge is reproduced with the radial flow plus trigger bias mechanism.  The composition of the \ridge is expected to be comparable to that in the inclusive spectra, as shown in \Fref{InclusiveRatioPlot}, because particles in the \ridge are from the bulk.  This is in agreement with Recombination because Recombination describes the inclusive spectra.   No calculation of the particle composition has been done in the Color Glass Condensate model, however, if it is able to describe the inclusive particle spectra, it should also describe the particle composition of the \ridge.
% 
% The trigger particle type dependence expected in any model incorporating hydrodynamical flow may provide a test for the radial flow plus trigger bias model.  As shown in \Fref{LowpTv2}, lighter particles are affected more by flow.  Hydrodynamical flow results in the movement of particles in response to a pressure gradient.  Assuming that the effective cross sections are comparable for both heavy and light particles, the forces on both are comparable, so the increase in velocity due to flow is greater for lighter quarks.  If the \ridge is produced by radial flow, a lighter parton will be correlated with a larger \ridge because the parton will experience a greater radial velocity boost, moving it towards the surface.  \Fref{RidgeNpartIDTrig} indicates that if there is a mass ordering, heavier hadron triggers are correlated with a greater \ridge.  If measurements of the \ridge for identified trigger particles can be done more precisely, it may be able to eliminate some of the proposed production mechanisms for the \ridge.
% 
\subsection{Summary}
The origin of the \ridge is not yet conclusive.  More quantitative theoretical calculations are needed in order to compare the data with the models.  Gluon bremsstrahlung models are likely excluded by the data because they are apparently in contradiction with the observations on the reaction plane dependence of the \ridge, the trigger particle type dependence of the \ridge, and the extent of the \ridge in \deta.  This does not mean that these mechanisms do not modify hadron fragmentation in heavy ion collisions, only that they are not likely to be the cause of the \ridge.  Indeed, gluon bremsstrahlung could cause the observed broadening of the \jlc in central \Au collisions at \sNNtwohundred.

The Momentum Kick Model is consistent with the data so far.  A calculation of the reaction plane dependence is crucial to determine if this model is in agreement with the data. Measurements which would test the prediction of a sharp drop of the \ridge with \deta are also crucial.

Models which rely on the radial flow plus trigger bias mechanism qualitatively agree with the data.  More precise calculations are needed.  Radial flow plus trigger bias would be the simplest mechanism to describe the \ridge since hydrodynamics describes the azimuthal asymmetry of particles produced in heavy ion collisions and radial flow has been observed in the inclusive spectra of particles.  If the radial flow plus trigger bias model is sufficient to explain the data, there is no need for more complicated, speculative mechanisms to produce the \ridge.  The significance of the \ridge in such a scenario would be that it would show up as a background in methods such as jet reconstruction and it is vital to understand how this background would affect reconstructed jets.  It also may be possible to learn something quantitative about radial flow from the \ridge; more theoretical developments along these lines would be necessary to determine if this is possible.
